% ---
% id: protocol-compiler/thm:chain-consumption
% kind: thm
% title: "Decided identities and named obligations of a consumed chain"
% summary: >-
%   Chain conditions are decided by finite computation; a verified pullback
%   identity reproduces the coarse-grained physical statistics and a verified
%   source identity maps the virtual source onto the specified one. Link
%   soundness and attack-set licences stay obligations this theorem does not
%   discharge.
% prominence: standard
% status: accepted
% uses: [def:chaincond, def:exact, thm:measurements-and-statistics, def:validator, thm:validator-soundness, def:target, def:compile, def:declaration, def:spec, transmitter/def:source-segment]
% mentions: [rem:coarse, rem:relabel, lem:weighted, def:weighted, def:banda, ass:atk, thm:gates, thm:band-a, reductions/thm:two-sided, rem:obligations-discharged, transmitter/rem:segment-scope]
% assumes: [A-LINK]
% refs:
%   - "tupkary2026decoy, Lemma 9.3, Eq. (146)"
%   - "tupkary2026decoy, Lemma 9.1, Eqs. (138)-(139)"
%   - "tupkary2026decoy, Rem. 9.2"
% provenance: newly-derived
% implements: certiqs_key.protocol_compiler.chain:check_chain
% lean:
% notes: "A-LINK is declared because the consistency reading of (iii) and (iii') is conditional on trace preservation of the supplied channel data, an occurrence of the assumption in the statement; A-ATK is a mention, spent at thm:gates (iii). reductions/thm:two-sided is a mention: (v') names it, and no step composes a source segment with a receiver segment. def:spec is a use: Step 3 reads the physical family as the POVM of its clause (b) for the clause of (iii) on the opposite direction."
% ---

Consuming a chain buys only what its record makes decidable. The theorem
states what each decided identity implies and lists what stays an obligation
of the chain's producer; where those obligations are discharged for the chains
the reductions module emits is Remark~\ref{rem:obligations-discharged}.
The attack-set licence that (v) names is Assumption~\ref{ass:atk}. It occurs
in no hypothesis of this theorem: it is spent at the gate that tests for its
record, Theorem~\ref{thm:gates}\,(iii), and is named here and not declared.

How the compiler applies the Bob-side weighting of (ii) after the chain is
Remark~\ref{rem:relabel}. Remark~\ref{rem:coarse} reads the clause of (iii)
on a coarse-graining in the opposite direction as the reason (C2) runs from
the physical alphabet onto the emitted one. That (A1)--(A3) do not decide whether the claim an obligation
record names applies to the descriptors of the link carrying that record,
as (v$'$) says, is
\xref{transmitter/rem:segment-scope}.

\begin{theorem}
\label{thm:chain-consumption}
Let a chain be presented to the compiler, with the
\tref[interface conditions on a chain]{interface conditions} of
Definition~\ref{def:chaincond}. Then:
\begin{enumerate}[label=(\roman*),leftmargin=2.2em]
\item (C1)--(C5) and (S1)--(S4) are decided by finite enumeration and finitely
many arithmetic operations on the emitted data, and a chain failing any of
them is refused. (X1) is not decided,
because it is not a condition on data: it is the type of the record;
\item (C3) is exactly the part of D(iii) concerning $\{F_v\}$, in both halves
--- typing against $\Smap$, and the per-block predicates and summation
condition of Definition~\ref{def:exact} --- hence a chain passing the interface
conditions yields an emitted $\{F_v\}$ for which Theorem~\ref{thm:measurements-and-statistics}
applies. Where the Bob-side weighting of Definition~\ref{def:weighted} is
applied after the chain (Definition~\ref{def:compile}), (C3) is decided on the family the chain delivered,
the emitted family is its weighting $N_{v,f}=P(f\,|\,t(v))F_v$, check V3
decides the emitted family itself (Theorem~\ref{thm:measurements-and-statistics}\,(i)), and
Lemma~\ref{lem:weighted}\,(i), (ii) is what says the weighting cannot break
what (C3) established; the identity of conjunct lists asserted here is between
(C3) and D(iii) on the chain's family;
\item \emph{(what a verified (C4) buys)} when the composite receiver channel
$\Lambda$ is given as data, (C4) is \emph{equivalent} to the statement that
measuring the emitted family $\{F_v\}$ after $\Lambda$ reproduces the
\emph{$\gamma$-coarse-grained} physical measurement statistics:
\begin{equation}\label{eq:pullbackstats}
  \tr[F_v\,\Lambda(\rho)]=\sum_{w\in\gamma^{-1}(v)}\tr[\Fph_w\,\rho]
\end{equation}
for every $v\in\Vset$ and every operator $\rho$ on the physical receiver space.
Consistency with completeness holds \emph{under one hypothesis this contract
names and does not verify}: if the supplied channel data are trace-preserving,
so that $\Lambda^{\dagger}$ is unital, then summing \eqref{eq:pullback} over
$v$ returns completeness of the physical family and nothing more, whereas the
same identity read with a coarse-graining in the opposite direction --- a
total surjection $\gamma':\Vset\to\Vph$, with
$\Lambda^{\dagger}(F_v)=\Fph_{\gamma'(v)}$ --- does not in general, forcing
every physical element covered more than once by $\gamma'$ to vanish. That
trace preservation is
Assumption~\ref{ass:link}, not a check: nothing here evaluates it;
\item \emph{(what a verified (S3) buys)} when the composite source channel
$\Xi$ is given as data, (S3) of
\xref{transmitter/def:source-segment} --- the source identity,
that definition's equation, decided by this module on the presented data
(Definition~\ref{def:chaincond}) --- says that the virtual source, pushed
through $\Xi$, is the specified source: $\Xi(\rvirt_x)=\rho_x$ for every $x$,
hence
\begin{equation}\label{eq:srcresid}
  (\idmap_X\otimes\Xi)\Big(\sum_{x\in\Xset}p_x\dyad{x}_X\otimes\rvirt_x\Big)
  =\sum_{x\in\Xset}p_x\dyad{x}_X\otimes\rho_x ,
\end{equation}
and, averaging, $\sum_xp_x\Xi(\rvirt_x)=\sum_xp_x\rho_x$: the
prepare-and-measure states of the emitted virtual description map onto the
specified ones, index by index, with the weights unchanged by (S2). The
compiler decides the source identity (S3) and nothing beyond it; and
\item the following are \emph{named and never discharged here}: that each link
of either segment is a sound reduction --- that its channel exists as claimed
for a link supplying only a theorem constructor, that a link supplying
its channel as data supplies a completely positive and trace-preserving map,
and that every reachable state after the reduction satisfies every transported
constraint entry --- and that any restricted \tref{attack set} claimed by a
link is licensed. These are obligations of the chain's producer, the
soundness of each link being Assumption~\ref{ass:link}. No claim of this contract discharges them for a
presented chain in general --- for the \tref{Band A} identity link, which
Definition~\ref{def:banda}\,(b) fixes with its channel supplied as data, its
outcome correspondence and its target POVM,
Theorem~\ref{thm:band-a}\,(v) discharges every clause of
Assumption~\ref{ass:link} by exhibition --- and
no proof below invokes them.
\end{enumerate}
\emph{The general format.} For a chain presented as
Definition~\ref{def:spec}\,(c$'$), with the interface conditions read as in the
general-format paragraph of Definition~\ref{def:chaincond}, (i)--(v) hold with
the following readings, of which the five above are the case of one
\tref{measured factor} and an empty \tref{source segment}.
\begin{enumerate}[label=(\roman*$'$),leftmargin=2.6em]
\item (C1)--(C5) are decided per measured factor and (S1)--(S4) as before, all
by finite enumeration and finitely many arithmetic operations, and a chain
failing any of them is refused.
\emph{And (A1)--(A3), the licensed-only rule and the declared-source-mode
rule are decided as well}: the first three by comparing finitely many declared
strings and one finite list against another, the fourth by reading $m$
declared tags, and the fifth by reading one declared slot. (X1) is again not
decided, being the type of the record.
\item (C3) is exactly the part of D(iii) concerning the declared measured
families, \emph{factor by factor}: for each $f$, (C3) read for $f$ and
D(iii)'s clauses on $\{F^{f}_{v_f}\}$ are the same list of conjuncts, so a
chain passing the interface conditions yields per-factor families for which
Theorem~\ref{thm:measurements-and-statistics}\,(i$'$) applies. \emph{That the product family is a POVM on
$\Qsys$ is on neither side of that pairing}, being a consequence of the
factors' conditions and no condition on data
(Definition~\ref{def:declaration}\,(DX3), Theorem~\ref{thm:validator-soundness}, general
format). The Bob-side weighting sentence applies per factor, to that factor's
delivered family.
\item \emph{(what a verified (C4) buys, per factor and jointly)} when every
link of segment $f$ supplies its channel as data,
\eqref{eq:pullback-factor} is equivalent to
$\tr[F^{f}_{v_f}\Lambda_f(\rho_f)]
=\sum_{w_f\in\gamma_f^{-1}(v_f)}\tr[F^{\mathrm{ph},f}_{w_f}\rho_f]$ for every
$v_f$ and every operator $\rho_f$ on factor $f$'s physical receiver space.
\emph{And the joint statement follows for every operator, product or not}: if
every segment supplies its channels as data then, with
$\Lambda\coloneqq\bigotimes_f\Lambda_f$, $\gamma\coloneqq\prod_f\gamma_f$ and
$F^{\mathrm{ph}}_w\coloneqq\bigotimes_fF^{\mathrm{ph},f}_{w_f}$,
\begin{equation}\label{eq:pullback-product}
  \Lambda^{\dagger}(F_v)=\sum_{w\in\gamma^{-1}(v)}F^{\mathrm{ph}}_w ,
  \qquad\text{hence}\qquad
  \tr[F_v\,\Lambda(\rho)]=\sum_{w\in\gamma^{-1}(v)}\tr[F^{\mathrm{ph}}_w\,\rho]
\end{equation}
for every $v\in\Vset$ and every operator $\rho$ on the product physical space.
\emph{Neither side of \eqref{eq:pullback-product} is formed or decided by
anything}: it is a consequence of the per-factor identities, and the objects in
it --- the product channel, the product physical family, the product target
family --- are data of no one. The consistency reading holds under the same
hypothesis read per factor: if every factor's supplied channel data are trace
preserving, so that each $\Lambda_f^{\dagger}$ is unital, then so is
$\Lambda^{\dagger}$ and summing \eqref{eq:pullback-product} over $v$ returns
completeness of the product physical family and nothing more. That trace
preservation is Assumption~\ref{ass:link}, per factor, and is evaluated
nowhere.
\item (S3) is as in (iv), with the split by kind of witness: what a
\emph{decided} (S3) buys is \eqref{eq:srcresid} and its averaged form, for the
data links' share of the composite; for a link supplying a theorem constructor
this module decides nothing and what its share buys is what the claim its
obligation record names concludes, which this contract neither states nor
asserts.
\item the list of obligations named and never discharged here gains three
items, each per the general format: that each link of \emph{each} receiver
segment is a sound reduction, factor by factor; that a non-empty source
segment composes with the receiver segments as a reduction, which is
\xref{reductions/thm:two-sided}, used by no step here; and that the claim an
obligation record names applies to that link's descriptors, which (A1)--(A3)
do \emph{not} decide. These are again the obligations of (v),
Assumption~\ref{ass:link} and the attack-set licence, and the emitting
module's own assertion; no claim of this contract discharges them, save as (v) says of
the \tref{Band A} identity link, and no proof below invokes them.
\end{enumerate}
\end{theorem}

\begin{proof}
\emph{Idea.} Items (i) and (i$'$) count what each condition decides. Items
(ii) and (ii$'$) compare two lists of conjuncts. Items (iii) and (iii$'$) are
the Hilbert--Schmidt duality between an adjoint identity and a statistics
identity, both exact, with trace preservation of the supplied channels as the
one hypothesis of the consistency reading. Items (iv) and (iv$'$) are the
linearity of a channel, and (v) and (v$'$) are statements about what this
contract asserts.

\emph{Step 1 ((i), decidability).} (C1) and (S1) compare finitely many
declared descriptors --- for (C1)'s last clause, the four components of a
sector map, the same combinatorial comparison V1 makes on the map of a
declared $\Qsys$-operator family and taken here on the descriptor, at (C1);
(C2) enumerates a finite table and checks totality and surjectivity by
counting preimages; (C3) is check V3 of Definition~\ref{def:validator}
restricted to $\{F_v\}$, namely $|\SecSet|(2|\Vset|+1)$ exact decisions on
blocks of size $d_s$ together with $|\SecSet|$ summation decisions, finite by
Theorem~\ref{thm:validator-soundness}\,(iv); (C4) decides $|\Vset|$ matrix identities on
matrices of the declared dimensions --- the adjoint of a channel given as data
(by Kraus operators or as a matrix acting on operators) is computed by finitely
many arithmetic operations, and each right-hand side of \eqref{eq:pullback} is
a finite sum over the tabulated fibre $\gamma^{-1}(v)$; (S2) compares two finite
index sets and $|\Xset|$ declared weights; (S3) decides $|\Xset|$ matrix
identities on matrices of the declared dimensions; (C5) and (S4) check the
presence of fields. (X1) is a statement about which of the two lists a link was
written into, which the record answers by construction.

\emph{Step 2 ((ii), the conjunct lists).} (ii) is the statement of (C3)
against D(iii), conjunct by conjunct. D(iii) asks two things of the emitted
$\{F_v\}$: that it be a family of block lists against $\Smap$, and that each
block family $\{F_v^{(s)}\}_v$ satisfy the predicates and the summation
condition of Definition~\ref{def:exact}. (C3) states the first as the type of
the object the last link delivers and the second as those same conditions, on
the emitted blocks. The two lists of conjuncts are identical, so (C3) is
D(iii) restricted to $\{F_v\}$ and Theorem~\ref{thm:measurements-and-statistics} applies to what the
chain delivered. Nothing in (ii) is inferred from the physical family, so a
chain that merely \emph{claims} a POVM property cannot pass one off as
checked: the predicates are decided on the matrices the compiler emits. The
block-diagonality half, by contrast, is a property the type carries and no
check confirms. An entry that the chain's producer has \emph{mis-assigned} to
a block is a well-typed object no predicate can catch. V1 refuses a
mis-assignment that also disagrees about the sectors; one that does not is a
defect of the chain's producer under Assumption~\ref{ass:link}, named there
and not discharged here.

\emph{Step 3 ((iii), what a verified (C4) buys).} The \tref{measurement channel} of
$\{F_v\}$ sends $\sigma$ to $\sum_v\tr[F_v\sigma]\dyad{v}$; the
$\gamma$-coarse-graining of the physical measurement channel sends $\rho$ to
$\sum_v\big(\sum_{w\in\gamma^{-1}(v)}\tr[\Fph_w\rho]\big)\dyad{v}$, the fibres
$\gamma^{-1}(v)$ partitioning $\Vph$ because $\gamma$ is a total surjection
(C2). Two maps into the same classical register agree on all inputs iff their
diagonal entries agree, i.e.\ iff \eqref{eq:pullbackstats} holds; and by the
definition of the adjoint with respect to the Hilbert--Schmidt inner product,
$\tr[F_v\Lambda(\rho)]=\tr[\Lambda^{\dagger}(F_v)\rho]$, so
\eqref{eq:pullbackstats} holds for all $\rho$ iff \eqref{eq:pullback} holds for
all $v$ --- the second direction because an operator whose trace against every
$\rho$ vanishes is the zero operator, the trace pairing being non-degenerate.
Both directions are available here because both sides are exact identities.
For the consistency claim, summing \eqref{eq:pullback} over $v$ gives
$\Lambda^{\dagger}\big(\sum_vF_v\big)=\sum_{w\in\Vph}\Fph_w$; with
$\sum_vF_v=\idon{\Qsys}$ from (C3) and $\Lambda^{\dagger}$ unital --- which is
where trace preservation of the supplied data is used, and it is
Assumption~\ref{ass:link} and not a check --- the left-hand side is
$\idon{}$, so what (C4) returns is $\sum_w\Fph_w=\idon{}$: completeness of the
physical family, a condition that family already satisfies, so (C4) adds no
restriction beyond the two completeness conditions.
The opposite direction, under the same hypothesis, is the identity
$\Lambda^{\dagger}(F_v)=\Fph_{\gamma'(v)}$ for every $v\in\Vset$, with
$\gamma':\Vset\to\Vph$ a total surjection. Summing it over $v$, the left-hand
side is again $\Lambda^{\dagger}\big(\sum_vF_v\big)=\idon{}$, by linearity of
$\Lambda^{\dagger}$, (C3) and unitality. On the right, the fibres
$\gamma'^{-1}(w)$ partition $\Vset$ because $\gamma'$ is total, so grouping
the terms by the value of $\gamma'$ gives
$\sum_{w\in\Vph}|\gamma'^{-1}(w)|\,\Fph_w=\idon{}$. The physical family is a
POVM (Definition~\ref{def:spec}\,(b)): $\sum_w\Fph_w=\idon{}$ and every
$\Fph_w\succeq0$. Subtracting,
$\sum_{w\in\Vph}\big(|\gamma'^{-1}(w)|-1\big)\Fph_w=0$, and every coefficient
is a non-negative integer because $\gamma'$ is surjective, so every term is
positive semidefinite. A vanishing sum of positive semidefinite operators has
every term zero: its trace is a sum of non-negative traces, so each term has
trace zero, and a positive semidefinite operator of trace zero is zero. Hence
$\Fph_w=0$ for every $w$ with $|\gamma'^{-1}(w)|\ge2$. That is a restriction
the physical family need not satisfy, so this reading does not in general
return completeness and nothing more. It does exactly when every physical
element covered more than once already vanishes, in particular when $\gamma'$
is a bijection: then every non-zero term of
$\sum_w|\gamma'^{-1}(w)|\,\Fph_w$ has coefficient one, and the summed
identity is $\sum_w\Fph_w=\idon{}$.
In the bijective case $\gamma$ is a bijection,
\eqref{eq:pullback} reads $\Lambda^{\dagger}(F_v)=\Fph_{\gamma^{-1}(v)}$, and
that is the published form of the condition,
\cite[Lemma~9.3, Eq.~(146)]{tupkary2026decoy}, read at the level of operators
--- there the two families are indexed by one alphabet.

\emph{Step 4 ((iv), what a verified (S3) buys).} By linearity,
$\idmap_X\otimes\Xi$ sends $\sum_xp_x\dyad{x}_X\otimes\rvirt_x$ to
$\sum_xp_x\dyad{x}_X\otimes\Xi(\rvirt_x)$; substituting
$\Xi(\rvirt_x)=\rho_x$ from the source identity (S3) gives
\eqref{eq:srcresid}, and tracing out $X$ gives its averaged form. The step is
linear in each $\rvirt_x$ and uses no property of the declared weights. That
the converse implication --- security of the virtual protocol transfers to the
real one --- requires more than the source identity is exactly the content of
Assumption~\ref{ass:link}: the published template is
\cite[Lemma~9.1, Eqs.~(138)--(139)]{tupkary2026decoy}, whose hypothesis is the
existence of such a $\Xi$ and whose conclusion this module does not assert.

\emph{Step 5 ((v), what is asserted).} (v) is a statement about what this
contract asserts. Its content is that Assumption~\ref{ass:link} and the
attack-set licence are named as obligations of the chain's producer, and
that no proof above invokes them: no step invokes Assumption~\ref{ass:link}
as an established obligation, its trace-preserving clause entering (iii) only
as the stated hypothesis of the consistency reading; no step invokes the
attack-set licence; and no step uses a theorem of the reductions
module. The discharge (v) names is Theorem~\ref{thm:band-a}'s, for the
Band-A identity link alone, and no step above rests on it. The
trace-preserving clause of Assumption~\ref{ass:link} is the hypothesis under
which (iii)'s consistency reading holds, so this theorem's statement is
conditional on it and declares it. The order rule (X1) is not an exception --- it is a
\emph{typing} property of the emitted record, and its warrant is the published
one: \cite[Rem.~9.2]{tupkary2026decoy} resolves the tension between source
maps, which are formulated for unrestricted attack sets, and squashing maps,
which restrict them, by applying the source-map lemma first, so that
restrictions arise only afterwards. The compiler enforces the order; it does
not claim that either lemma applies to a given device.

\emph{The general format.} The steps are those above, read per factor.

\emph{Step 1 in the general format ((i$'$)).} Each clause of (i) is the
clause of (i) applied to one segment, the data being per factor and finite;
the refusal naming the factor is the record's own field. For (A1): the
registry an artefact emits is a finite list of triples of declared strings,
the obligation record a triple of declared strings, and the comparison is
entry by entry, including the two components of the assumption list; the
ledger membership test is a finite list against a finite list. For (A3): the
provenance test reads which module emitted the artefact the segment is
presented in, a declared field. The licensed-only rule reads $m$ tags,
terminating and exact. The declared-source-mode rule reads one declared slot,
whether the \tref{marginal slot} holds $\nomarg$, beside whether the source
segment is empty, which the record answers by its shape. None of the five is
an operator predicate and none is decided at a window.

\emph{Step 2 in the general format ((ii$'$)).} For each $f$, D(iii) asks of
$\{F^{f}_{v_f}\}$ that it be a family of block lists against factor $f$'s map
and that each block family satisfy the predicates and the summation condition
of Definition~\ref{def:exact}, and (C3) read for $f$ states the type and those
same conditions on that factor's emitted blocks; the two lists of conjuncts
are identical, factor by factor. What Step 2 says of the physical family and
of the block-diagonality half applies per factor unchanged, the
mis-assignment exposure now being one per factor and V1's per-factor map
comparison what catches an assignment that also disagrees about that factor's
sectors.

\emph{Step 3 in the general format ((iii$'$)).} The per-factor equivalence is
(iii)'s argument with $\Qsys_f$, $\Lambda_f$, $\gamma_f$ and factor $f$'s two
families in place of $\Qsys$, $\Lambda$, $\gamma$ and the two families: the
fibres $\gamma_f^{-1}(v_f)$ partition factor $f$'s physical alphabet because
$\gamma_f$ is a total surjection (C2), the adjoint identity is the
Hilbert--Schmidt pairing, and both directions are available because both
sides are exact identities. For \eqref{eq:pullback-product}: the adjoint of a
tensor product of maps is the tensor product of their adjoints, so
$\Lambda^{\dagger}(F_v)=\bigotimes_f\Lambda_f^{\dagger}(F^{f}_{v_f})$ by
\eqref{eq:productfamily}, this is
$\bigotimes_f\sum_{w_f\in\gamma_f^{-1}(v_f)}F^{\mathrm{ph},f}_{w_f}$ by
\eqref{eq:pullback-factor}, and expanding the product of sums by
multilinearity gives
$\sum_{w\in\gamma^{-1}(v)}\bigotimes_fF^{\mathrm{ph},f}_{w_f}$, the index set
being exactly $\gamma^{-1}(v)=\prod_f\gamma_f^{-1}(v_f)$. The statistics form
follows for \emph{every} operator $\rho$ on the product physical space, product
or not, because $\tr[F_v\Lambda(\rho)]=\tr[\Lambda^{\dagger}(F_v)\rho]$ is
linear in $\rho$ and \eqref{eq:pullback-product} is an identity of operators.
For the consistency reading, a tensor product of unital maps is unital, so
summing over $v$ and using $\sum_{v_f}F^{f}_{v_f}=\idon{\Qsys_f}$ per factor
(C3) gives $\sum_{w}F^{\mathrm{ph}}_w=\idon{}$, which the product physical
family already satisfies; nothing is added and nothing is evaluated.

\emph{Steps 4 and 5 in the general format ((iv$'$) and (v$'$)).} The first is
(iv)'s computation, which is linear in each $\rvirt_x$ and reads only the
identity, together with the split by kind of witness made in (S3) of
\xref{transmitter/def:source-segment}, restated here and not re-derived. The second is
again a statement about what this contract asserts: no step above invokes
Assumption~\ref{ass:link} as an established obligation, per factor or
otherwise, or the attack-set licence, and none uses a theorem of the
reductions module, and the three items named are named and not used --- in particular no
step above composes a source segment with a receiver segment, which is why the
composition rule appears in the list rather than in a hypothesis. What is
declared here is unchanged: the trace-preserving clause of
Assumption~\ref{ass:link} occurs in (iii$'$)'s consistency reading as in
(iii)'s, and the attack-set licence is named here and spent at the gate
(Theorem~\ref{thm:gates}\,(iii)).
\end{proof}
